% =============================================================================
% Kapitel 6: Diskussion
% =============================================================================
\chapter{Diskussion}
\label{ch:diskussion}

\section{Aussagekraft der Strukturbestimmung}

Die \evax{}-Analyse liefert ein atomistisches Strukturmodell, das die
experimentellen EXAFS-Spektren aller beteiligten Kanten simultan
reproduziert.
Daraus lassen sich partielle Koordinationszahlen, Bindungslängen und
die chemische Nahordnung ableiten -- Größen, die mit konventionellen
Fitmethoden bei Mehrstoffsystemen nur unter starken Annahmen
zugänglich sind.

Gleichzeitig ist zu bedenken, dass das Modell nicht eindeutig ist:
Die endliche Superzelle (256~Atome) stellt eine Stichprobe der
realen Struktur dar, und verschiedene Atomkonfigurationen können
ähnliche Residuen ergeben.
Die Verwendung des evolutionären Algorithmus mit 16~parallelen States
mildert dieses Problem, beseitigt es aber nicht grundsätzlich.
Die Streuung der Ergebnisse über die States liefert ein Maß für die
statistische Unsicherheit des Modells.

\section{Rolle der Parameteroptimierung}

Die Parameterstudie (Abschnitt~\ref{sec:ergebnisse-paramscan}) zeigt,
dass die Wahl des Vergleichsraums und die Berücksichtigung von
Mehrfachstreuung den größten Einfluss auf die Fitqualität haben.
Der Wavelet-Raum erweist sich als überlegen, was mit seiner
gleichzeitigen $k$- und $R$-Auflösung konsistent ist
(vgl.\ Abschnitt~\ref{sec:wavelet}).

Bemerkenswert ist, dass die strukturellen Ergebnisse
(Koordinationszahlen, SRO-Parameter) über die verschiedenen
Parametersätze hinweg \TODO{robust / leicht variabel} sind.
Dies spricht dafür, dass die physikalisch relevante Information
tatsächlich in den Daten enthalten ist und nicht nur ein Artefakt der
Parameterwahl darstellt.

\section{Unsicherheiten und systematische Fehler}

\subsection{$\chi(k)$-Extraktion}

Die Normierung der $\chi(k)$-Daten hat direkten Einfluss auf die
Amplitude des EXAFS-Signals und damit auf die Koordinationszahlen.
Bei der vorliegenden Multi-Edge-Messung der MEA-Probe musste die
Normierung für jede Kante separat durchgeführt werden
(Abschnitt~\ref{sec:datenvorverarbeitung}).
Die Erfahrung zeigt, dass eine um \SI{10}{\percent} fehlerhafte
Normierung zu einer entsprechenden Skalierung der
Koordinationszahlen führt.

\subsection{Korrelation von $S_0^2$ und $N$}

Der Amplitudenreduktionsfaktor $S_0^2$ und die Koordinationszahl $N$
sind physikalisch verschieden, treten in der EXAFS-Gleichung
(Gl.~\eqref{eq:exafs}) jedoch als Produkt $S_0^2 \cdot N$ auf.
Die automatische $S_0^2$-Bestimmung durch \evax{} liefert Werte, die
mit Literaturangaben für metallische Systeme konsistent sind
($S_0^2 \approx 0{,}6$--$0{,}9$).
Eine unabhängige Kalibrierung an einer Referenzprobe würde die
Zuverlässigkeit weiter erhöhen.

\subsection{Modellannahmen}

Das Strukturmodell setzt ein periodisch fortgesetztes Kristallgitter
voraus.
Für die Bulk-Legierung (Probe~1) ist diese Annahme gerechtfertigt.
Bei Nanopartikeln (Probe~2) hingegen weichen die Koordinationszahlen
an der Oberfläche deutlich von den Bulk-Werten ab, was zu einer
systematischen Unterschätzung der mittleren Koordinationszahl
führen kann.
\evax{} bietet die Möglichkeit, endliche Cluster statt periodischer
Zellen zu verwenden, was bei einer weiterführenden Analyse zu
berücksichtigen wäre.

\section{Einordnung in die Literatur}

\TODO{Nach Vorliegen der Ergebnisse: Vergleich mit veröffentlichten
Werten für CoNiCu-MEAs (z.\,B.\ Gitterparameter aus XRD, SRO aus
Monte-Carlo-Simulationen) und CoPt-Nanopartikel (L1$_0$-Ordnung,
Partikelgröße aus TEM).}

Für das CoNiCu-System ist bekannt, dass alle drei reinen Metalle
oberhalb von ca.\ \SI{700}{\kelvin} in der FCC-Phase vorliegen und
eine lückenlose Mischkristallreihe bilden \cite{George2019}.
Die Frage nach der chemischen Nahordnung -- ob bestimmte
Nachbarpaare statistisch bevorzugt oder vermieden werden -- ist
hingegen experimentell wenig untersucht und stellt daher den
eigentlichen Informationsgewinn der vorliegenden Analyse dar.

\section{Vergleich mit konventionellen Fitmethoden}

Im Vergleich zu einem konventionellen EXAFS-Fit in \textsc{Artemis}
bietet die \evax{}-Analyse zwei wesentliche Vorteile:
\begin{enumerate}[nosep]
  \item Die Zahl der effektiven Fitparameter skaliert nicht mit der
        Zahl der Streupfade, da die Atompositionen selbst die einzigen
        freien Parameter sind.
  \item Das Modell liefert unmittelbar dreidimensionale Strukturen,
        aus denen sich höhere Korrelationsfunktionen (z.\,B.\ die
        Dreikörperverteilung) extrahieren lassen.
\end{enumerate}

Der Nachteil liegt in der erheblich höheren Rechenzeit und der
Abhängigkeit vom Startmodell: Während ein Shell-Fit in Minuten
abgeschlossen ist, benötigt ein \evax{}-Lauf mit 256~Atomen und
drei Kanten mehrere Stunden.
Für Routineanalysen bleibt der konventionelle Ansatz daher das Mittel
der Wahl; für komplexe Mehrstoffsysteme, bei denen die
Parameterkorrelation zu groß wird, stellt \evax{} eine leistungsfähige
Alternative dar.
